\chapter{प्रस्तावना}

इस किताब की महत्वाकांक्षा सीधी-सादी है: जीव विज्ञान की एक ही सुसंगत किताब,
जो किंडरगार्टन से लेकर स्नातक की डिग्री के आख़िर तक सब कुछ समेटे, और जो
दुनिया में कहीं भी पढ़ने वालों के लिए लिखी गई हो।

इसकी शैली छोटी और कड़ी है: परिभाषाओं, उदाहरणों, प्रतिज्ञप्तियों, तरीक़ों
और चित्रों से बना पाठ, जहाँ भी उस स्तर पर तर्क समझ में आ सके वहाँ ध्यान
से रखा गया तर्क, और उसके बाद सोच-समझकर चुने गए अभ्यास। जो नतीजे पूरी
व्युत्पत्ति के बिना कहे गए हैं उन्हें साफ़-साफ़ स्वीकृत लिखा गया है। सारे
अभ्यासों के पूरे हल किताब के अंत में इकट्ठे हैं। और जीव विज्ञान के अध्याय
समीकरणों से ज़्यादा चित्रों को तरजीह देते हैं।

स्कूल का हर साल इस किताब का एक भाग है, और हर भाग अध्यायों में बँटा है।
साल जितना पहले का, पढ़ने वाला उतना ही छोटा: यानी ज़्यादा चित्र, ज़्यादा
खोलकर बताए गए क़दम, और ज़्यादा नरम अभ्यास।

\section*{इस किताब को कैसे पढ़ें}

कथनों की गिनती हर अध्याय के भीतर होती है और वे रंगों से पहचाने जाते हैं:
परिभाषाएँ नीली, प्रमेय लाल, प्रतिज्ञप्तियाँ तथा उनके रिश्तेदार नारंगी, और
तरीक़े --- यानी आम कामों के छोटे नुस्ख़े --- हरे। अभ्यास सीधे इस्तेमाल
वाले ($\star$) से लेकर कठिन समस्याओं ($\star\star\star$) तक दर्जों में
बँटे हैं।

\section*{योगदान}

इस किताब का स्रोत एक सार्वजनिक git भंडार में रहता है:
\begin{center}
\url{https://github.com/bvirrion/one-biology-book}
\end{center}
योगदान --- नए अध्याय, सुधार, बेहतर उपपत्तियाँ, और अतिरिक्त अभ्यास ---
सबका स्वागत है; इसके लिए भंडार की जड़ में रखी \texttt{CONTRIBUTING.md}
फ़ाइल देखें।
